\documentclass[12pt]{article}
\usepackage[utf8]{inputenc}
\usepackage[T1]{fontenc}
\usepackage{hyperref}
\title{Tubar Language Resources for K-12 Curriculum Development}
\author{Compiled by Claude Sonnet 5 (claude-sonnet-5) for Andr\'as Kornai}
\date{}
\begin{document}
\maketitle
\section{Introduction}
Tubar (ISO 639-3 tbu) is an extinct Uto-Aztecan language of southern Chihuahua, Mexico. Its placement within the family is debated (a ``curious blend'' of Taracahitan and Tepiman features). It was recorded around 1893 by Carl Hartman and Carl Lumholtz, when fewer than 50 speakers remained; the data were analysed and published only later (Lionnet 1978). Very little is available online.
\section{Primary Educational Resources}
No primers, textbooks or school materials were found in the searches.
\section{Linguistic and Reference Materials}
\begin{itemize}
\item Wikipedia, Tubar language: short sketch, including a Lord's Prayer sample. \url{https://en.wikipedia.org/wiki/Tubar_language}
\item Haugen, Uto-Aztecan slides (2014), discussing the classification of Tubar. \url{https://www.eva.mpg.de/fileadmin/content_files/linguistics/conferences/2014-Mesoamerican/Haugen_Uto-Aztecan_Slides.pdf}
\item Numeral data (one to five only). \url{https://lingweb.eva.mpg.de/channumerals/Tubar.htm}
\item Lionnet (1978), \emph{El idioma tubar y los tubares}, and Stubbs (2000), comparative phonology, are cited in the sources above; I did not see them online, so no URLs are given.
\end{itemize}
\section{Religious and Community Resources}
None found.
\section{Digital and Audio Resources}
None found. No recordings exist, as there are no known speakers.
\section{Curriculum Development Context}
Tubar is extinct and documentation is minimal (a few word lists, numerals 1--5, short texts). A curriculum could at most be a heritage-awareness unit built on the historical record; it cannot use the language for instruction. Any real work would require retrieving Lionnet (1978) and the original Hartman/Lumholtz records.
\end{document}